\section{Introdução}
\label{sec:intro}

Modelos de linguagem de grande escala (LLMs) vêm sendo incorporados a sistemas
multiagente para apoiar análise financeira: interpretação de notícias, leitura
de indicadores, construção de teses de investimento e tomada de decisão em
ambientes de trading. O volume de propostas cresce, mas uma lacuna
metodológica permanece. Grande parte dos trabalhos atribui ganhos à presença
de múltiplos agentes ou ao uso de debate sem comparar os resultados contra
baselines fortes de agente único submetidos ao mesmo orçamento de
computação. Fica difícil, então, distinguir o que vem da arquitetura
multiagente do que vem, simplesmente, de um modelo mais capaz acompanhado
de prompts mais bem feitos.

A questão pesa especialmente no contexto financeiro. Uma resposta que soa
convincente pode não significar maior personalização ao investidor, melhor
cobertura dos fatores de risco da carteira ou menor taxa de erro factual
sobre os ativos analisados. Soma-se a isso o fato de que a maior parte dos
sistemas existentes prioriza recomendação de ativos ou decisão de compra e
venda, deixando em segundo plano um problema mais frequente na prática:
diagnosticar uma carteira já constituída em relação ao perfil de risco, ao
horizonte e ao objetivo do usuário. Diagnosticar não é o mesmo que
recomendar. O objetivo deixa de ser propor movimentações sobre os ativos e
passa a ser explicar o estado atual da carteira, apontar fatores de
atenção e sinalizar quando a evidência disponível é insuficiente.

Este trabalho propõe uma arquitetura multiagente baseada em LLMs para
diagnóstico personalizado de carteiras financeiras. A saída do sistema é
uma classificação discreta da carteira frente ao perfil informado, com
justificativa textual, fatores positivos, fatores negativos e nível de
confiança. Não há recomendações de compra, venda ou rebalanceamento, e o
sistema não assume papel de aconselhamento financeiro. Para sustentar
afirmações sobre o impacto da arquitetura, a proposta é avaliada contra um
conjunto de baselines que inclui agente único ingênuo, agente único com
raciocínio estruturado e self-consistency, e multiagente sem debate, sob
mesma entrada e mesmo formato de saída.

As principais contribuições deste trabalho são:

\begin{itemize}
  \item uma arquitetura multiagente para diagnóstico personalizado de
  carteiras, organizada em agentes especialistas (técnico, fundamentalista,
  sentimento, macroeconômico e de risco), um par adversarial Bull/Bear e um
  moderador que produz a classificação final;
  \item um mecanismo de debate adversarial Bull/Bear com síntese por
  moderador e suporte a roteamento heterogêneo de modelos por papel, que
  permite isolar o efeito da arquitetura do efeito do modelo subjacente;
  \item uma metodologia de avaliação que compara o sistema proposto contra
  baselines fortes de agente único, self-consistency e multiagente sem
  debate, usando entrada congelada e saída estruturada para garantir
  comparações controladas;
  \item um pipeline de rastreabilidade que registra, por execução, chamadas
  realizadas, latência, tokens, erros de parsing e respostas estruturadas,
  permitindo reproduzir e auditar tanto a arquitetura proposta quanto os
  baselines.
\end{itemize}
